% GENERATED FILE. Edit canonical lessons, then run npm run generate:books.
% canonical-source-hash: fnv1a64:6bbd0c0d4e29fd48
% canonical-lessons: GE-C383-behoerde, GE-C383-antrag, GE-C383-beantragen, GE-C383-formular, GE-C383-unterschrift

\chapter{Bei der Behörde}
\label{ch:ge-bei-der-behorde}
\hlchaptermodality{383}

\begin{chapteropening}
\textbf{By the end of this chapter:} \emph{I can deal with an authority: make an application, fill in a form and sign it.}

Complete the last lesson of chapter 383: I can deal with an authority: make an application, fill in a form and sign it.
\end{chapteropening}

\section[die Behörde]{die Behörde — the authority, the public office}

\label{lesson:GE-C383-behoerde}



\begin{quote}
Before the new one: say the German for to break, then the German for the violation.
\end{quote}

\subsection*{You'll want to know: die Behörde}
\textbf{die Behörde} — \textquotedblleft{}the authority, the public office\textquotedblright{}.

\begin{quote}
Für den neuen Ausweis muss ich zur Behörde gehen.
\end{quote}

\begin{quote}
\emph{For the new identity card I have to go to the authorities.}
\end{quote}

Plural: \emph{die Behörden}. In Switzerland people often say \emph{das Amt}.

\begin{cousinweb}
\emph{Behörde} is usually traced to an old verb \emph{behören} \textquotedblleft{}to belong to\textquotedblright{}, of the family of \emph{hören} \textquotedblleft{}to hear\textquotedblright{} and \emph{gehören} \textquotedblleft{}to belong\textquotedblright{}: the office a matter belongs to. EtymDB has no history for it, so the account is hedged.
\end{cousinweb}

\begin{grammarlens}[title={What you've built}]
One more word for telling, explaining and describing at B1.
\end{grammarlens}

\subsection*{Your turn}
\begin{itemize}
\raggedright
  \item \emph{Say it:} \emph{die Behörde}
  \item \emph{Say it:} \emph{die Behörde}, once more
  \item \emph{Say it:} a sentence of your own with \emph{die Behörde}
  \item \emph{Recall:} say the German for to break, then the German for the violation, then say \emph{die Behörde} again
\end{itemize}

\begin{tcolorbox}[breakable,colback=teal!4,colframe=teal!35!black,title={Before you move on}]
What does die Behörde mean? (\textquotedblleft{}The authority, the public office\textquotedblright{}.) Say it once more.
\end{tcolorbox}
\section[der Antrag]{der Antrag — the application; the motion}

\label{lesson:GE-C383-antrag}



\begin{quote}
Before the new one: say the German for the violation, then the German for the authority.
\end{quote}

\subsection*{You'll want to know: der Antrag}
\textbf{der Antrag} — \textquotedblleft{}the application; the motion\textquotedblright{}.

\begin{quote}
Ich habe einen Antrag auf Wohngeld gestellt.
\end{quote}

\begin{quote}
\emph{I have applied for housing benefit.}
\end{quote}

Plural: \emph{die Anträge}. \emph{Einen Antrag stellen} — to make an application. \emph{Der Heiratsantrag} — the marriage proposal.

\begin{cousinweb}
\emph{Antrag} is inherited from Middle High German \emph{antrac} (EtymDB), of the family of \emph{antragen} \textquotedblleft{}to offer\textquotedblright{}, built on \emph{tragen} \textquotedblleft{}to carry\textquotedblright{}: you carry a request to the office.
\end{cousinweb}

\begin{grammarlens}[title={What you've built}]
One more word for telling, explaining and describing at B1.
\end{grammarlens}

\subsection*{Your turn}
\begin{itemize}
\raggedright
  \item \emph{Say it:} \emph{der Antrag}
  \item \emph{Say it:} \emph{der Antrag}, once more
  \item \emph{Say it:} a sentence of your own with \emph{der Antrag}
  \item \emph{Recall:} say the German for the violation, then the German for the authority, then say \emph{der Antrag} again
\end{itemize}

\begin{tcolorbox}[breakable,colback=teal!4,colframe=teal!35!black,title={Before you move on}]
What does der Antrag mean? (\textquotedblleft{}The application; the motion\textquotedblright{}.) Say it once more.
\end{tcolorbox}
\section[beantragen]{beantragen — to apply for}

\label{lesson:GE-C383-beantragen}



\begin{quote}
Before the new one: say the German for the authority, then the German for the application.
\end{quote}

\subsection*{You'll want to know: beantragen}
\textbf{beantragen} — \textquotedblleft{}to apply for\textquotedblright{}.

\begin{quote}
Sie hat ein Visum beantragt.
\end{quote}

\begin{quote}
\emph{She has applied for a visa.}
\end{quote}

Inseparable: \emph{sie beantragt, sie hat beantragt}. It takes a direct object, with no \emph{für}.

\begin{cousinweb}
\emph{beantragen} is built with \emph{be-} on \emph{Antrag} \textquotedblleft{}application\textquotedblright{}, from Middle High German \emph{antrac} (EtymDB), of the family of \emph{tragen} \textquotedblleft{}to carry\textquotedblright{}.
\end{cousinweb}

\begin{grammarlens}[title={What you've built}]
One more word for telling, explaining and describing at B1.
\end{grammarlens}

\subsection*{Your turn}
\begin{itemize}
\raggedright
  \item \emph{Say it:} \emph{beantragen}
  \item \emph{Say it:} \emph{beantragen}, once more
  \item \emph{Say it:} a sentence of your own with \emph{beantragen}
  \item \emph{Recall:} say the German for the authority, then the German for the application, then say \emph{beantragen} again
\end{itemize}

\begin{tcolorbox}[breakable,colback=teal!4,colframe=teal!35!black,title={Before you move on}]
What does beantragen mean? (\textquotedblleft{}To apply for\textquotedblright{}.) Say it once more.
\end{tcolorbox}
\section[das Formular]{das Formular — the form (to fill in)}

\label{lesson:GE-C383-formular}



\begin{quote}
Before the new one: say the German for the application, then the German for to apply for.
\end{quote}

\subsection*{You'll want to know: das Formular}
\textbf{das Formular} — \textquotedblleft{}the form (to fill in)\textquotedblright{}.

\begin{quote}
Das Formular hat zwölf Seiten.
\end{quote}

\begin{quote}
\emph{The form has twelve pages.}
\end{quote}

Plural: \emph{die Formulare}. \emph{Ein Formular ausfüllen} — to fill in a form.

\begin{cousinweb}
\emph{Formular} is usually traced to Latin \emph{formulārius}, which Lewis \& Short give for a lawyer who drew up legal forms, from \emph{formula} \textquotedblleft{}a small form, a pattern\textquotedblright{}, the diminutive of \emph{forma} \textquotedblleft{}shape\textquotedblright{} (Lewis \& Short).
\end{cousinweb}

\begin{grammarlens}[title={What you've built}]
One more word for telling, explaining and describing at B1.
\end{grammarlens}

\subsection*{Your turn}
\begin{itemize}
\raggedright
  \item \emph{Say it:} \emph{das Formular}
  \item \emph{Say it:} \emph{das Formular}, once more
  \item \emph{Say it:} a sentence of your own with \emph{das Formular}
  \item \emph{Recall:} say the German for the application, then the German for to apply for, then say \emph{das Formular} again
\end{itemize}

\begin{tcolorbox}[breakable,colback=teal!4,colframe=teal!35!black,title={Before you move on}]
What does das Formular mean? (\textquotedblleft{}The form (to fill in)\textquotedblright{}.) Say it once more.
\end{tcolorbox}
\section[die Unterschrift]{die Unterschrift — the signature}

\label{lesson:GE-C383-unterschrift}



\begin{quote}
Before the new one: say the German for to apply for, then the German for the form.
\end{quote}

\subsection*{You'll want to know: die Unterschrift}
\textbf{die Unterschrift} — \textquotedblleft{}the signature\textquotedblright{}.

\begin{quote}
Auf dem Formular fehlt noch Ihre Unterschrift.
\end{quote}

\begin{quote}
\emph{Your signature is still missing on the form.}
\end{quote}

Plural: \emph{die Unterschriften}. From \emph{unterschreiben} — to sign.

\begin{cousinweb}
\emph{Unterschrift} joins \emph{unter} \textquotedblleft{}under\textquotedblright{} to \emph{Schrift} \textquotedblleft{}writing\textquotedblright{}, from Old High German \emph{scrift} and Latin \emph{scrīptum} (EtymDB): the name written under the text.
\end{cousinweb}

\begin{grammarlens}[title={What you've built}]
One more word for telling, explaining and describing at B1.
\end{grammarlens}

\subsection*{Your turn}
\begin{itemize}
\raggedright
  \item \emph{Say it:} \emph{die Unterschrift}
  \item \emph{Say it:} \emph{die Unterschrift}, once more
  \item \emph{Say it:} a sentence of your own with \emph{die Unterschrift}
  \item \emph{Recall:} say the German for to apply for, then the German for the form, then say \emph{die Unterschrift} again
\end{itemize}

\begin{tcolorbox}[breakable,colback=teal!4,colframe=teal!35!black,title={Before you move on}]
What does die Unterschrift mean? (\textquotedblleft{}The signature\textquotedblright{}.) Say it once more.
\end{tcolorbox}
